\chapter{Results}
\label{chap: Results}

% Opening: 1-2 paragraphs. State what is shown and in what order.
% Specifically: the chapter leads with the headline hybrid-loop demonstration
% (§4.2), then decomposes the result by validating each component in turn
% (§4.3 AE, §4.4 NODE, §4.5 OOD), and closes with the generalisation /
% transfer-learning study (§4.6) and an aggregate ablation summary (§4.7).
% Recap of the test case (cylinder, Re = 2500, 2D, n = 2^9 grid) to anchor
% all the numbers that follow. One sentence on hardware (DelftBlue, single
% NVIDIA GPU, Float32) so the wall-times in §4.2 have a referent.

% -----------------------------------------------------------------------
\section{Experimental Setup and Conventions}
\label{sect: 4_setup}



% Short. Lifts the relevant subset of §3.1/§3.8 into a single place the reader
% can refer back to without flipping chapters:
%   - Domain, grid, Re, U/L, boundary treatment (1 paragraph).
%   - Train / val / test temporal split (figure or sentence).
%   - Hardware, precision, seed.
%   - Hyperparameters used for the "baseline" hybrid run: latent_dim = 16,
%     group_size = 20, continuity_term = 200, λ_div, λ_curl, n_switch = 150,
%     Δt* = 0.35, max_retrain_flags = 3. Everything that follows is this run
%     unless stated.

% -----------------------------------------------------------------------
\section{Hybrid Acceleration of the Cylinder Flow}
\label{sect: 4_hybrid}

% The headline section. This is the section you can draft now from existing
% plt_combined.png runs.

\subsection{Force Coefficients and Wall Time}
\label{subsect: 4_forces_walltime}
% Figure: plt_combined.png-style 3-panel: C_D / C_L vs t (hybrid + reference),
% mode log (warmup / train / hybrid / retrain colour bands), wall-time bar.
% Numbers: mean |C_D - C_D_ref|, RMS(C_L - C_L_ref), overall_speedup.

\subsection{Mean Flow and Reynolds-Stress Statistics}
\label{subsect: 4_meanflow_rst}
% Figure: rst_comp_plot + plt_meanflow side by side.
% Quantitative: spatial L2 errors on <u>, <v>, <u'u'>, <v'v'>, <u'v'>.

\subsection{Physical Consistency of the Hybrid State}
\label{subsect: 4_physical_consistency}
% Divergence residual ||∇·û|| over time, max & mean.
% Snapshot vorticity comparison hybrid vs reference at 3 instants.
% Important framing: the hybrid trajectory is one realisation of a chaotic
% system, so pointwise matching is the wrong test; statistical matching is
% the right one. State this explicitly.

% -----------------------------------------------------------------------
\section{Autoencoder Reconstruction Quality}
\label{sect: 4_ae_results}


\subsection{Reconstruction Error on Held-Out Snapshots}
\label{subsect: 4_ae_recon}
% Figure: side-by-side ground truth / reconstruction / |error| for a handful
% of test snapshots. From visualize_reconstructions().
% Table: MAE, correlation coefficient per channel, divergence residual, curl
% residual on the test split.

\subsection{Effect of Latent Dimension}
\label{subsect: 4_ae_latent}
% From your hpc_tune sweep. Plot recon_mae and divergence vs N_z for
% N_z ∈ {8, 16, 32, 64} or whatever you swept. Justify N_z = 16 as the
% knee of the curve.

\subsection{Effect of Physics-Informed Loss Weights}
\label{subsect: 4_ae_phys_loss}
% λ_div, λ_curl sweep. Show that without them the reconstruction is still
% sharp but divergence-free is violated; with them, recon MAE rises slightly
% but div/curl drop by an order of magnitude. This is the empirical
% justification for the loss design in §3.3.2.

\subsection{Training Dynamics}
\label{subsect: 4_ae_training}
% From plot_losses(): train / val / test loss curves, correlation curves.
% One figure, brief discussion of convergence behaviour and any overfitting.

% -----------------------------------------------------------------------
\section{Neural ODE Latent Rollouts}
\label{sect: 4_node_results}

\subsection{Rollout Accuracy Against Encoded Ground Truth}
\label{subsect: 4_node_rollout}
% Per-dimension plot z_i(t) for ground truth vs NODE rollout, train region
% and test region clearly marked. The "extrapolation" plot. State the
% horizon at which rollout error crosses a threshold.

\subsection{Effect of Multiple-Shooting Parameters}
\label{subsect: 4_node_ms}
% Sweep over group_size ∈ {10, 15, 20, 30} and continuity_term ∈
% {100, 200, 400} (your sweep_node_retrain.jl grid). Plot rollout error
% (and possibly training stability) vs each. Justify group_size = 20,
% continuity_term = 200 from the data.


% -----------------------------------------------------------------------
\section{OOD Detector Calibration and Behaviour}
\label{sect: 4_ood_results}

\subsection{Threshold Calibration on Training Latents}
\label{subsect: 4_ood_threshold}
% Histogram of KNN distances on training latents, with the 99th-percentile
% threshold τ marked.


\subsection{Sensitivity to k and threshold}
\label{subsect: 4_ood_sensitivity}
% Brief: how the trigger rate during the hybrid loop changes with k and τ.

% -----------------------------------------------------------------------
\section{Generalisation}
\label{sect: 4_transfer}

% This is where TL1 / TL2 live. From the existing checkpoints
% TL1_E500_..., TL2_E300_..., and sweep_node_retrain.jl.

\subsection{Transfer to a Finer Simulation Mesh}
\label{subsect: 4_tl_re}
% Hybrid run at Re ≠ 2500 using TL1-fine-tuned AE. Show what fraction of
% training cost is recovered relative to training from scratch.


% -----------------------------------------------------------------------
\section{Summary of Findings}
\label{sect: 4_summary}

% Short. One paragraph each on (i) end-to-end acceleration achieved,
% (ii) which design choices the ablations validated, (iii) which honest
% limitations are visible in the results (most likely: chaotic divergence
% in pointwise comparisons, fixed geometry, BC overhead, single Re without
% TL). The committee will look for this last paragraph.